% 第五章 系统核心模块实现

\chapter{系统核心模块实现}

\newcommand{\chapteruiscreenshotplaceholder}[2]{
    \begin{figure}[htbp]
        \centering
        \IfFileExists{figures/screenshots/#1}{
            \includegraphics[width=0.9\textwidth]{figures/screenshots/#1}
        }{
            \fbox{
                \parbox[c][5.2cm][c]{0.86\textwidth}{
                    \centering
                    界面截图待补充\\[6pt]
                    对应界面：#2
                }
            }
        }
        \caption{#2}
        \label{fig:#1}
    \end{figure}
}

本章在第四章总体设计的基础上，按照系统实现职责重新组织FreeFlow的核心模块。与原先按功能点逐项展开的写法不同，本文将实现内容归并为客户端以及智能合约与服务器两个部分。其中，客户端侧重点说明用户可见界面、交互流程和本地播放体验；智能合约与服务器侧重点说明身份认证、IPFS存储、链上发布、访问控制和链下资源投影等支撑能力。



\section{客户端}

客户端是FreeFlow面向用户的统一交互入口，采用Electron与React实现，既承担创作者发布作品时的表单编辑、素材管理和链上交易发起，也承担普通用户检索资源、查看详情、购买访问权、播放音频和发表评论等操作。为了使论文展示更贴近系统实际运行效果，本节以界面组织方式说明客户端实现。



\subsection{用户资料引导与SIWE登录界面}

用户首次进入系统时，客户端首先展示资料引导界面。该界面承担三项任务：选择或创建本地Profile、连接钱包并切换到目标链、完成SIWE（Sign-In with Ethereum）登录。界面布局上，左侧用于展示已有钱包Profile列表，右侧用于显示当前钱包连接状态、链网络状态以及登录操作按钮，使用户能够在同一页面完成身份建立与业务会话初始化。

在交互实现上，客户端先读取本地保存的Profile列表，并根据当前运行时上下文恢复最近一次使用的钱包身份；当用户点击连接钱包后，系统调用钱包连接组件拉起外部钱包，并在必要时引导用户切换到Sepolia测试网络。钱包连接成功后，客户端向后端申请nonce，构造SIWE消息，请求钱包签名，并将签名结果提交后端验签。若验证通过，客户端获得Web2.5会话信息，并把当前钱包地址与本地Profile建立绑定关系，从而保证“本地身份、链上地址、服务器会话”三者一致。

从界面反馈角度看，该页面需要持续展示连接进度、网络错误、签名取消和登录成功等状态。这样一来，钱包登录不再只是底层能力调用，而是被转化为用户能够理解和确认的逐步流程界面，适合在论文中作为系统启动入口展示。

\chapteruiscreenshotplaceholder{profile_guide.png}{用户资料引导与钱包登录界面}

\subsection{设置与动态主题模块}

\subsection{创作者工作台与发行编辑界面}

创作者工作台是客户端中最具代表性的复合界面，也是创作者完成链上音乐发布的主要入口。该界面以“导航栏 + 工作区 + 状态提示”的方式组织功能，支持发行项目列表、草稿创建、作品信息编辑、素材状态展示、元数据上传、链上发布和收益查看等操作。为了降低创作者对链上流程的理解门槛，工作台没有直接暴露复杂的合约细节，而是将发布过程拆解为多个连续步骤，并以状态标签驱动界面切换。

从界面结构上看，工作台顶部区域用于展示当前模块标题与摘要信息，左侧导航栏用于在“发布作品”“元数据编辑”“收益中心”等主视图之间切换，中部工作区则根据当前阶段渲染对应的表单、上传面板或活动记录。这样做的目的是将“内容编辑”“链上操作”“结果查看”拆分为清晰的视觉区域，避免创作者在单页中同时面对过多操作控件。

\begin{table}[htbp]
    \linespread{1.45}
    \zihao{5}
    \songti
    \centering
    \caption{创作者工作台主要界面区域说明}\label{tab:creator_workshop_regions}
    \begin{tabular}{|l|p{4.0cm}|p{6.1cm}|}
        \hline
        界面区域 & 主要内容 & 作用说明 \\ \hline
        顶部信息区 & 当前模块标题、摘要信息、用户状态 & 帮助创作者快速确认所处业务阶段，并提供统一的页面语义 \\ \hline
        左侧导航区 & 发布作品、元数据编辑、收益中心 & 对工作台中的不同工作面板进行切换，保持功能入口稳定 \\ \hline
        草稿列表区 & 已创建发行项目、状态标签、更新时间 & 支持创作者在多个作品草稿之间切换与继续编辑 \\ \hline
        编辑工作区 & 标题、描述、价格、封面、音频、歌词等表单 & 承载发行信息填写、素材上传和metadata生成等核心输入操作 \\ \hline
        状态与操作区 & 上传按钮、发布按钮、错误提示、链上结果 & 展示当前步骤是否完成，并限制无效操作顺序 \\ \hline
    \end{tabular}
\end{table}

发行状态主要包括草稿（\texttt{DRAFT}）、素材待完善（\texttt{ASSETS\_PENDING}）、素材已上传（\texttt{ASSETS\_UPLOADED}）、元数据已上传（\texttt{METADATA\_UPLOADED}）、发布中（\texttt{PUBLISHING}）、已发布（\texttt{PUBLISHED}）和失败（\texttt{FAILED}）。每个状态对应不同的操作权限和界面提示。例如，当音频或封面尚未上传完成时，发布按钮保持禁用；当交易发送后，界面则突出显示交易处理中状态，避免用户重复点击或误以为操作失效。

发行状态流转如图\ref{fig:release_state_flow}所示。工作台通过状态字段限制用户操作顺序，避免在素材未上传、metadata未生成或链上交易未确认时进入后续阶段。失败状态会保留错误信息，便于用户修复输入或网络问题后重新执行对应步骤。

\chapteruiscreenshotplaceholder{creators_workshop.png}{创作者工作台主界面}

\chapteruiscreenshotplaceholder{publish_form.png}{发行信息编辑与素材上传界面}

\begin{figure}[htbp]
    \centering
    \includegraphics[width=0.95\textwidth]{figures/out/release_state_flow.png}
    \caption{创作者发行状态流转图}
    \label{fig:release_state_flow}
\end{figure}

\subsection{链上资源搜索与详情入口界面}

虽然链上资源排序算法将在第六章专门讨论，但从用户体验角度看，客户端仍需要将搜索结果组织为可阅读、可比较、可跳转的界面。系统在搜索结果页中以列表方式展示封面、标题、创作者、价格、访问模式和链上身份摘要，并为每条结果提供进入详情页的统一入口。这样做一方面便于用户在多个链上资源之间快速浏览和筛选，另一方面也为后续的购买、加入链上音乐库和评论等行为提供稳定的导航路径。

该界面在展示策略上突出两类信息：一类是用户最关心的消费决策信息，例如作品标题、创作者名称、是否公开、价格和封面；另一类是链上可信信息，例如合约地址摘要、Token ID或资源唯一键。前者服务于可读性和审美展示，后者服务于可追溯性与可信验证。通过这种双层信息组织方式，搜索结果页既保持了音乐应用界面的直观性，也体现了链上资源的技术特征。

\chapteruiscreenshotplaceholder{search_results.png}{链上资源搜索结果界面}

\subsection{链上音乐详情、购买与链上音乐库界面}

用户从搜索结果进入作品详情页后，客户端会围绕“资源展示、访问状态、购买操作、链上凭证确认”四部分组织页面内容。详情页优先展示封面、标题、创作者、专辑或描述信息，并同时提供价格、访问模式、创作者地址、metadata链接和发布交易链接等链上字段，保证用户在视觉浏览之外仍能访问作品的可验证信息。

对于公开作品，详情页会直接显示“加入链上音乐库”入口；对于付费作品，界面则突出价格信息和购买按钮，并在按钮区域提示当前钱包是否具备访问权。购买流程启动后，客户端调用\texttt{PlatformHub}合约的\texttt{buyAccess}方法，等待交易确认后，再次读取链上访问状态以及ERC-1155余额结果，用于更新“已购买”“可播放”“可加入音乐库”等界面状态。该设计遵循“链上结果优先”的原则，即后端索引负责展示和检索，真正的访问资格仍由链上状态决定。

链上音乐库用于承接公开或已购买资源的长期管理。用户将资源加入链上音乐库后，客户端会把作品写入本地统一歌曲记录中，并附带\texttt{resourceKey}、链ID、合约地址、\texttt{tokenId}、metadata CID和音频CID等扩展字段，使链上资源能够复用本地播放器、歌单和缓存逻辑。这样，链上音乐的使用体验在界面层面就能够接近普通本地音乐。

\chapteruiscreenshotplaceholder{track_detail.png}{链上音乐详情与购买界面}

\subsection{歌单浏览与播放队列管理界面}

除链上资源消费外，FreeFlow仍然保留了完整的播放器组织能力。歌单界面用于承载用户日常的曲目浏览和批量播放操作，其页面通常由封面头图、歌单标题、平台标识、播放全部按钮以及曲目列表构成。用户既可以在本地歌单中管理普通音乐，也可以在链上音乐库歌单中管理已经加入库的FreeFlow资源，因此该界面实际上承担了“本地内容管理”和“链上资源消费承接”两类任务。

在具体实现上，歌单页面支持按照标题、艺术家和添加时间排序，并支持关键字过滤。用户点击“播放全部”后，客户端会将当前歌单曲目整体替换到播放队列中，再由播放器按照当前播放模式开始播放。这种设计使歌单不只是静态列表，而是与队列控制、播放器状态和链上音乐库记录直接联动，从而形成统一的消费入口。

\chapteruiscreenshotplaceholder{playlist.png}{歌单浏览与批量播放界面}

播放队列界面进一步展示系统在播放组织上的动态能力。该界面通常位于右侧抽屉区域，包含“队列”和“最近播放”两个标签页，并分别展示当前播放曲目、接下来待播放曲目以及历史播放记录。对于接下来待播放的条目，系统允许用户通过双击立即切换播放，并支持清空队列或清空最近播放记录。这样，用户不仅能够查看当前正在消费的内容，也能够理解后续播放顺序，从而获得比传统“上一首/下一首”按钮更强的控制感。

从论文展示角度看，播放队列界面能够很好体现播放器内部状态的可视化实现。它将\texttt{PlayQueue}中的当前曲目、剩余曲目和历史曲目以可见列表形式展示出来，使播放状态不再只是控制器内部数据，而是可以被用户直接观察和调整的界面能力。这一部分对于说明系统具备完整播放器实现，而非仅有链上购买展示页面，尤其重要。

\chapteruiscreenshotplaceholder{play_queue.png}{播放队列与最近播放界面}

\subsection{播放控制、歌词联动与评论交互界面}

播放器界面承担实际消费体验，是客户端实现中最接近传统音乐应用的部分。底部播放器栏负责展示当前封面、标题和艺术家信息，并提供上一首、播放/暂停、下一首、播放模式切换、队列展开、歌词跳转和音量调节等核心控件。进度条部分支持多种皮肤样式，但在功能上都支持拖动定位、缓冲区显示和当前时间同步更新，从而保证播放控制具有较好的可视反馈。

歌词界面则体现了播放器在“播放实现”上的联动能力。系统在歌词页左侧展示封面或唱片动画，右侧展示滚动歌词区，并支持原文、音译和译文等不同歌词类型切换。播放过程中，客户端根据当前播放时间定位活跃歌词行，并自动滚动到对应位置；用户也可以直接点击某一行歌词，将播放进度跳转到对应时间点。这样，歌词页不只是附属信息显示，而是成为可参与播放定位的交互界面。

\chapteruiscreenshotplaceholder{lyric.png}{歌词联动与播放控制界面}

FreeFlow没有让界面直接访问远程网关音频，而是通过本地音频代理统一处理播放请求。播放器请求本地代理接口后，代理服务优先读取本地文件，其次读取磁盘缓存，最后再根据资源解析服务得到的IPFS网关地址拉取音频流。代理服务支持Range请求，因此用户可以拖动进度条进行分段播放，避免在播放长音频时重复下载完整文件。

音频代理与缓存流程如图\ref{fig:audio_proxy_cache_flow}所示。该流程将本地文件、缓存文件和远程音频源统一封装为本地代理响应，降低界面层直接处理跨域、Range和缓存细节的复杂度，也使链上音乐、本地音乐和外部来源音乐在播放器中获得较一致的使用体验。

\begin{figure}[htbp]
    \centering
    \includegraphics[width=0.95\textwidth]{figures/out/audio_proxy_cache_flow.png}
    \caption{音频代理与缓存流程图}
    \label{fig:audio_proxy_cache_flow}
\end{figure}

在互动功能方面，评论界面由作品评论列表、评论输入区和登录提示区组成。用户登录后可以针对特定发行资源发表评论，未登录时界面则提示先完成钱包登录。在实际交互上，歌词页和详情页都可以继续引导用户进入评论区域，因此评论功能并不是孤立存在，而是被嵌入到“浏览---播放---反馈”的消费路径中。这样一来，客户端既展示了播放器的消费能力，也展示了作品围绕播放行为产生的社区互动能力。


\section{智能合约与服务器}

相较于客户端负责用户可见的交互过程，智能合约与服务器共同承担系统可信业务与链下协同能力。前者负责访问凭证、销售规则和收益分账等不可篡改逻辑，后者负责SIWE认证、Pinata上传代理、发行草稿管理、资源检索、购买记录与评论服务等链下支撑模块。二者协同后，系统能够在保持链上可信性的同时提供较完整的业务可用性。

\subsection{认证、资料与会话服务实现}

系统首先通过后端认证服务完成钱包身份验证。后端生成nonce后，桌面端构造SIWE消息并请求钱包签名，后端对签名消息进行解析和验签。验证通过后，系统创建或更新用户、钱包身份和登录会话，并将会话令牌同步到桌面端业务客户端。

该模块的核心安全边界在后端。桌面端只负责发起钱包连接、展示签名内容和保存会话令牌，后端负责校验nonce是否过期或已消费、签名地址是否与请求地址一致、chainId是否符合要求，以及domain、uri和issuedAt等SIWE字段是否有效。验证成功后，会话以令牌哈希形式保存到\texttt{auth\_sessions}表，桌面端后续访问受保护接口时通过Authorization头携带会话令牌。这样既避免了客户端单独决定身份状态，也为评论、发行草稿和收益查询等接口提供统一的身份基础。

\subsection{素材上传与metadata生成服务实现}

创作者选择音频和封面后，客户端将文件提交给后端Pinata上传接口。后端读取文件内容，调用Pinata API上传，并把返回CID、文件大小、MIME类型和Pinata ID记录到数据库，用于后续资源追踪与错误排查。此处由服务器代为上传而非客户端直连Pinata，一方面便于统一保存上传记录，另一方面也避免在客户端暴露敏感配置。

元数据JSON文件由客户端根据作品信息、音频CID、封面CID、访问模式、价格、链上合约地址和歌词信息生成，并通过后端存储服务上传。元数据上传成功后，系统将\texttt{metadataStorageObjectId}和\texttt{metadataDocument}回写到发行记录中，使后续链上发布能够直接引用metadata URI。

\begin{table}[htbp]
    \linespread{1.5}
    \zihao{5}
    \songti
    \centering
    \caption{元数据JSON主要字段说明}\label{tab:metadata_json_fields}
    \begin{tabular}{|l|p{9.5cm}|}
        \hline
        字段             & 说明                                                        \\ \hline
        name            & 作品标题，用于链上\texttt{tokenURI}对应元数据的基础展示                           \\ \hline
        description     & 作品描述，来源于创作者工作台编辑内容                                      \\ \hline
        image           & 封面网关地址或IPFS引用，用于详情页和音乐库封面展示                          \\ \hline
        animation\_url  & 音频网关地址或IPFS引用，用于播放器解析远程音频资源                           \\ \hline
        attributes      & 结构化扩展属性，可包含艺术家、专辑、流派、访问模式、试听秒数和链上合约信息等                \\ \hline
        external\_url   & 作品外部链接，可指向后端资源详情或区块浏览器页面                               \\ \hline
        freeflow        & 系统扩展字段，可保存audioCid、coverCid、metadataCid、访问模式和来源信息等 \\ \hline
    \end{tabular}
\end{table}

\subsection{链上发布、购买授权与收益分账实现}

元数据准备完成后，创作者可调用\texttt{PlatformHub}合约的\texttt{publishTrack}方法发布作品。该方法内部创建收益分账合约、生成音乐访问凭证\texttt{tokenId}，并保存作品的销售配置。交易确认后，客户端解析\texttt{TrackPublished}事件，获取\texttt{tokenId}、创作者地址和分账合约地址，并将交易哈希、区块号和\texttt{tokenId}回写后端。这样，链上合约承担的是权威状态记录，后端承担的是链下索引补全与展示投影。

购买流程同样由\texttt{PlatformHub}负责。对于需要付费访问的作品，用户调用\texttt{buyAccess}方法完成支付；合约校验销售状态和支付金额后，调用\texttt{MusicAccess1155}铸造访问凭证，并将创作者收益转入对应的\texttt{RoyaltySplitter}分账合约。此后，客户端和后端都不单独声称用户“已经拥有访问权”，而是通过\texttt{hasAccess}和ERC-1155余额结果来确认链上授权状态。

收益分账由\texttt{RoyaltySplitter}合约完成。创作者可配置多个收款地址及其份额，购买收益进入分账合约后，各收款人可按份额提取。相较于将收入在单个地址中集中管理，独立的分账合约更便于体现协作创作场景下的收益规则，也使平台能够以合约形式固化收入分配逻辑。

\subsection{资源索引、访问投影与评论服务实现}

服务器数据库以\texttt{creator\_releases}为核心构建链上资源索引，用于支撑资源搜索、详情解析和资源唯一键定位。用户在客户端打开FreeFlow链上音乐详情页后，系统根据\texttt{resourceKey}解析后端资源索引，并读取链上销售配置。后端提供的是“可检索、可展示、可联动”的链下视图，而不是对链上状态的替代。

在购买场景中，系统会将购买结果写入\texttt{purchases}表，主要目的不是替代链上记录，而是为购买历史展示、资源热度统计和重复写入防护提供链下依据。当用户再次进入详情页时，系统可以结合链上状态与数据库投影更快地还原消费上下文，但最终访问资格仍由合约结果决定。

评论模块则完全由后端评论服务驱动。用户登录后可以针对特定发行资源发表评论，后端保存评论内容、用户标识和时间信息；客户端再据此展示评论列表、作者信息和发布时间。评论虽然不属于链上可信状态，但它直接增强了作品传播和用户互动能力，是连接资源消费与社区交流的重要链下模块。


\section{本章小结}

本章从客户端与智能合约、服务器两个视角介绍了FreeFlow的核心实现。客户端部分重点展示了钱包登录、创作者工作台、链上资源详情、播放器和评论等可见界面；智能合约与服务器部分重点说明了SIWE认证、IPFS素材上传、metadata生成、链上发布、购买授权、收益分账和评论索引等支撑机制。二者共同构成了从创作者发布到用户检索、购买、播放和互动的完整业务闭环。
